% ============================================================
%  CUADERNILLO DE ESTUDIO — CLASE 11
%  Seguridad de la Información y de Sistemas
%  Lic. en Gestión de Negocios Digitales — UCA
%  Compilar con: pdflatex (dos pasadas para el índice)
% ============================================================

\documentclass[12pt,a4paper]{article}

% ---------- Paquetes ----------
\usepackage[utf8]{inputenc}
\usepackage[spanish]{babel}
\usepackage[T1]{fontenc}
\usepackage{lmodern}
\usepackage[margin=2.2cm,headheight=14pt]{geometry}
\usepackage{amsmath}
\usepackage{amssymb}
\usepackage{enumitem}
\usepackage{booktabs}
\usepackage{array}
\usepackage{tabularx}
\usepackage[table]{xcolor}
\usepackage{tcolorbox}
\usepackage{graphicx}
\usepackage{fancyhdr}
\usepackage{titlesec}
\usepackage{hyperref}
\usepackage{multicol}
\usepackage{pifont}
\usepackage{wasysym}
\usepackage{tikz}
\usetikzlibrary{positioning,arrows.meta,shapes.geometric}

% ---------- Colores ----------
\definecolor{azulUCA}{HTML}{003366}
\definecolor{doradoUCA}{HTML}{B8860B}
\definecolor{grisClaro}{HTML}{F2F2F2}
\definecolor{rojoAlerta}{HTML}{C0392B}
\definecolor{verdeOK}{HTML}{27AE60}
\definecolor{azulCielo}{HTML}{2980B9}
\definecolor{naranjaAmenaza}{HTML}{D35400}
\definecolor{violetaParcial}{HTML}{6C3483}
\definecolor{azulGobierno}{HTML}{1B4F72}
\definecolor{verdePSA}{HTML}{117A65}

% ---------- tcolorbox estilos ----------
\tcbuselibrary{skins,breakable}

\newtcolorbox{cajaTitulo}[1][]{
  colback=azulUCA, colframe=azulUCA, coltext=white,
  fonttitle=\Large\bfseries, arc=4mm, #1
}

\newtcolorbox{cajaDefinicion}{
  colback=azulCielo!10, colframe=azulCielo,
  fonttitle=\bfseries, title=Definici\'on, arc=3mm, breakable
}

\newtcolorbox{cajaGobierno}{
  colback=azulGobierno!8, colframe=azulGobierno,
  fonttitle=\bfseries\color{azulGobierno},
  title={\ding{211}\; Gobierno de la Seguridad}, arc=3mm, breakable
}

\newtcolorbox{cajaPSA}{
  colback=verdePSA!8, colframe=verdePSA,
  fonttitle=\bfseries\color{verdePSA},
  title={\ding{45}\; Pol\'itica de Seguridad Aceptable}, arc=3mm, breakable
}

\newtcolorbox{cajaParcial}{
  colback=violetaParcial!8, colframe=violetaParcial,
  fonttitle=\bfseries\color{violetaParcial},
  title={\ding{115}\; PARCIAL 2}, arc=3mm, breakable
}

\newtcolorbox{cajaActividad}{
  colback=grisClaro, colframe=azulUCA,
  fonttitle=\bfseries\color{azulUCA},
  title={\ding{45}\; Actividad Pr\'actica}, arc=3mm, breakable
}

\newtcolorbox{cajaNota}{
  colback=doradoUCA!10, colframe=doradoUCA,
  fonttitle=\bfseries, title=Concepto clave, arc=3mm, breakable
}

% ---------- Header / Footer ----------
\pagestyle{fancy}
\fancyhf{}
\fancyhead[L]{\small\textcolor{azulUCA}{Seguridad de la Informaci\'on y de Sistemas}}
\fancyhead[R]{\small\textcolor{azulUCA}{Cuadernillo --- Clase 11}}
\fancyfoot[C]{\thepage}
\fancyfoot[L]{\small\textcolor{gray}{UCA}}
\renewcommand{\headrulewidth}{0.4pt}
\renewcommand{\footrulewidth}{0.4pt}

% ---------- Títulos de sección ----------
\titleformat{\section}
  {\color{azulUCA}\Large\bfseries}
  {\thesection.}{0.6em}{}
\titleformat{\subsection}
  {\color{azulCielo}\large\bfseries}
  {\thesubsection}{0.6em}{}

% ---------- Enlaces ----------
\hypersetup{colorlinks=true, linkcolor=azulUCA, urlcolor=azulCielo}

% ---------- Checkbox ----------
\newcommand{\checkboxVacio}{$\square$\;}
\newcommand{\checkMark}{\ding{51}}

% ============================================================
\begin{document}

% ============================================================
%  PORTADA
% ============================================================
\begin{titlepage}
\centering
\vspace*{1cm}

\begin{cajaTitulo}[width=\textwidth, halign=center]
  SEGURIDAD DE LA INFORMACI\'ON Y DE SISTEMAS\\[4pt]
  {\large Licenciatura en Gesti\'on de Negocios Digitales --- UCA}
\end{cajaTitulo}

\vspace{1.2cm}

{\Huge\bfseries\color{azulUCA} CUADERNILLO DE ESTUDIO}\\[8pt]
{\LARGE\color{doradoUCA} CLASE 11}\\[10pt]
{\Large\itshape ``Gobierno de la seguridad:\\[2pt] Pol\'itica de Seguridad Aceptable''}\\[6pt]
{\large\color{violetaParcial}\ding{115}\; Incluye PARCIAL 2}\\[1.5cm]

\begin{tabular}{rl}
\textbf{Profesor:}   & Mg.\ Ing.\ Rodrigo Atilio Elgueta \\[4pt]
\textbf{Unidad:}     & Unidad 4 \\[4pt]
\end{tabular}

\vfill

\begin{tcolorbox}[colback=grisClaro, colframe=azulUCA, arc=3mm, width=0.9\textwidth]
  \centering
  {\bfseries Datos del/la estudiante}\\[10pt]
  Nombre y apellido: \underline{\hspace{5cm}}\\[10pt]
  Grupo N.\textdegree: \underline{\hspace{5cm}}\\[10pt]
  Negocio Digital (TP): \underline{\hspace{5cm}}\\[10pt]
  Fecha: \underline{\hspace{5cm}}
\end{tcolorbox}

\vspace{1cm}
{\small\textcolor{gray}{Material de uso exclusivo para la cursada.}}
\end{titlepage}

% ============================================================
%  ÍNDICE
% ============================================================
\tableofcontents
\newpage

% ============================================================
%  SECCIÓN 1 — INTRODUCCIÓN
% ============================================================
\section{De la t\'ecnica al gobierno: el tablero de la seguridad}

En las unidades anteriores aprendimos a proteger sistemas (Unidad 1 y 2), a cumplir la ley (Unidad 3) y a concientizar personas (Clase 10). Pero todo eso puede quedar como acciones aisladas si la organizaci\'on no tiene \textbf{reglas claras del juego}.

\begin{cajaDefinicion}
El \textbf{Gobierno de la Seguridad} es el conjunto de pol\'iticas, normas, procedimientos, roles y responsabilidades que una organizaci\'on establece para dirigir, controlar y mejorar continuamente su seguridad de la informaci\'on.
\end{cajaDefinicion}

\begin{cajaNota}[title={La diferencia fundamental}]
La \textbf{seguridad t\'ecnica} dice \emph{c\'omo} se protegen los sistemas (firewalls, cifrado, MFA). El \textbf{gobierno de la seguridad} dice \emph{qui\'en decide qu\'e se protege, por qu\'e, y qu\'e pasa si no se cumple}. Sin gobierno, la seguridad es un conjunto de herramientas sin estrategia.
\end{cajaNota}

\newpage
% ============================================================
%  SECCIÓN 2 — PIRÁMIDE DOCUMENTAL
% ============================================================
\section{La pir\'amide documental: pol\'iticas, normas y procedimientos}

El gobierno de la seguridad se materializa en documentos organizados jer\'arquicamente. Cada nivel tiene un prop\'osito distinto.

\begin{center}
\begin{tikzpicture}[
  box/.style={rectangle, draw=azulGobierno, align=center, rounded corners=3pt},
]
  % Pirámide desde arriba (más estratégico) a abajo (más operativo)
  \node[box, fill=azulGobierno!60, text=white, minimum width=3cm, minimum height=1.4cm, font=\small\bfseries] (pol)
    {POL\'ITICAS\\{\scriptsize ¿QU\'E y POR QU\'E?}};

  \node[box, fill=azulGobierno!40, text=white, minimum width=6cm, minimum height=1.4cm, font=\small\bfseries, below=0.4cm of pol] (nor)
    {NORMAS\\{\scriptsize Reglas espec\'ificas y obligatorias}};

  \node[box, fill=azulGobierno!25, minimum width=9cm, minimum height=1.4cm, font=\small\bfseries, below=0.4cm of nor] (pro)
    {PROCEDIMIENTOS\\{\scriptsize Pasos detallados, c\'omo se ejecuta}};

  \node[box, fill=grisClaro, draw=gray, minimum width=11.5cm, minimum height=1.2cm, font=\small, below=0.4cm of pro] (gui)
    {GU\'IAS / INSTRUCTIVOS / PLANTILLAS (opcional)};

  \draw[-Stealth, thick, azulGobierno] (pol.south) -- (nor.north);
  \draw[-Stealth, thick, azulGobierno] (nor.south) -- (pro.north);
  \draw[-Stealth, thick, gray] (pro.south) -- (gui.north);
\end{tikzpicture}
\end{center}

\renewcommand{\arraystretch}{1.5}
\begin{tabularx}{\textwidth}{|>{\bfseries\color{azulGobierno}}p{2.8cm}|X|X|}
\hline
\rowcolor{grisClaro}
\textbf{Nivel} & \textbf{¿Qu\'e es?} & \textbf{Ejemplo} \\
\hline
\textbf{Pol\'itica} & Declaraci\'on de intenci\'on de la direcci\'on. Establece el ``qu\'e'' y el ``por qu\'e''. Es estable en el tiempo. & ``Todos los datos de clientes se cifran en reposo y tr\'ansito.'' \\
\hline
\textbf{Norma} & Regla obligatoria y espec\'ifica que desarrolla la pol\'itica. & ``Se usa AES-256 para datos en reposo; TLS 1.2 o superior para datos en tr\'ansito.'' \\
\hline
\textbf{Procedimiento} & Pasos detallados para cumplir la norma. Es operativo. & ``Paso 1: generar clave en el KMS. Paso 2: configurar el bucket S3 con SSE-KMS. Paso 3: auditar trimestralmente.'' \\
\hline
\end{tabularx}

\begin{cajaGobierno}[title={Por qu\'e importa el orden}]
Si una empresa solo tiene \textbf{procedimientos} sin pol\'iticas, nadie sabe \emph{por qu\'e} se hacen las cosas y todo depende de personas espec\'ificas. Si solo tiene \textbf{pol\'iticas} sin procedimientos, las buenas intenciones no se ejecutan. El equilibrio entre los tres niveles es lo que hace que la seguridad sea \textbf{sostenible}.
\end{cajaGobierno}

\newpage
% ============================================================
%  SECCIÓN 3 — PSA
% ============================================================
\section{La Pol\'itica de Seguridad Aceptable (PSA)}

\begin{cajaDefinicion}
La \textbf{Pol\'itica de Seguridad Aceptable} (en ingl\'es, Acceptable Use Policy, AUP) es el documento que establece \textbf{qu\'e pueden y qu\'e no pueden hacer los empleados y usuarios} con los recursos tecnol\'ogicos de la organizaci\'on (equipos, red, internet, correo, aplicaciones, datos).
\end{cajaDefinicion}

\subsection{¿Por qu\'e es tan importante?}
\begin{itemize}[leftmargin=2em]
    \item \textbf{Legal:} da respaldo para aplicar sanciones ante incumplimientos (un empleado que usa la notebook de la empresa para minar criptomonedas, por ejemplo).
    \item \textbf{Claridad:} elimina zonas grises; todo el mundo sabe qu\'e est\'a permitido.
    \item \textbf{Cultura:} instala la seguridad como parte del contrato laboral, no como una molestia externa.
\end{itemize}

\subsection{Estructura m\'inima de una PSA}

Una PSA profesional suele tener estos bloques:

\renewcommand{\arraystretch}{1.5}
\begin{tabularx}{\textwidth}{|>{\bfseries\color{verdePSA}}p{3.5cm}|X|}
\hline
\rowcolor{grisClaro}
\textbf{Secci\'on} & \textbf{¿Qu\'e contiene?} \\
\hline
1. Prop\'osito y alcance & ¿Para qu\'e existe la pol\'itica? ¿A qui\'en aplica? (empleados, contractors, terceros). \\
\hline
2. Uso aceptable de los recursos & Qu\'e usos est\'an permitidos (trabajo, uso personal razonable) y cu\'ales prohibidos. \\
\hline
3. Contrase\~nas y accesos & Requisitos m\'inimos de contrase\~nas, uso de gestores, MFA, prohibici\'on de compartir credenciales. \\
\hline
4. Correo e internet & Uso aceptable del correo corporativo y la conexi\'on a internet; sitios prohibidos; manejo de adjuntos. \\
\hline
5. Dispositivos & Uso de equipos corporativos, BYOD (``tra\'e tu propio dispositivo''), p\'erdida o robo. \\
\hline
6. Informaci\'on y datos confidenciales & Clasificaci\'on de datos, manejo, prohibici\'on de fugas, redes sociales. \\
\hline
7. Incidentes & Obligaci\'on de reportar incidentes de seguridad. \\
\hline
8. Cumplimiento y sanciones & Qu\'e pasa si no se cumple (desde advertencia hasta despido o acciones legales). \\
\hline
9. Firmas y vigencia & Firma del empleado, firma del responsable, fecha de vigencia y revisi\'on. \\
\hline
\end{tabularx}

\begin{cajaPSA}[title={Regla de oro de una buena PSA}]
Una PSA debe ser \textbf{corta, clara y espec\'ifica}. Si es un documento de 40 p\'aginas en lenguaje legal, nadie la lee y no sirve. Las mejores PSA caben en 2-4 p\'aginas y se firman al ingresar a la empresa.
\end{cajaPSA}

\newpage
% ============================================================
%  SECCIÓN 4 — ROLES Y RESPONSABILIDADES
% ============================================================
\section{Roles y responsabilidades en seguridad}

El gobierno tambi\'en define \textbf{qui\'en es responsable de qu\'e}. Los roles m\'as comunes en una organizaci\'on digital son:

\renewcommand{\arraystretch}{1.5}
\begin{tabularx}{\textwidth}{|>{\bfseries\color{azulGobierno}}p{3.2cm}|X|}
\hline
\rowcolor{grisClaro}
\textbf{Rol} & \textbf{Responsabilidad principal} \\
\hline
\textbf{Direcci\'on / CEO} & Aprueba la pol\'itica y asigna presupuesto. Es el responsable \'ultimo. \\
\hline
\textbf{CISO} (Chief Information Security Officer) & Lidera la estrategia de seguridad, reporta a la direcci\'on. \\
\hline
\textbf{DPO} (Data Protection Officer) & Vela por el cumplimiento de la normativa de datos personales (Ley 25.326 / GDPR). \\
\hline
\textbf{Due\~no del activo (Asset Owner)} & Responsable de un activo espec\'ifico (ej. el Gerente Comercial es due\~no de la base de clientes). \\
\hline
\textbf{\'Area de TI / Seguridad} & Implementa controles t\'ecnicos, monitorea, responde incidentes. \\
\hline
\textbf{Todos los empleados} & Cumplen la PSA, reportan incidentes, protegen sus credenciales. \\
\hline
\end{tabularx}

\begin{cajaNota}[title={El error frecuente en las pymes]}
En empresas chicas no suele haber un CISO ni un DPO dedicados. Pero \textbf{los roles deben existir aunque los ocupe la misma persona}. Lo importante es que alguien sea formalmente responsable de la seguridad, de los datos personales y del cumplimiento. Si ``todos son responsables'', nadie lo es.
\end{cajaNota}

\newpage
% ============================================================
%  SECCIÓN 5 — PARCIAL 2: CASO Y RÚBRICA
% ============================================================
\section{PARCIAL 2: caso, formato y r\'ubrica}

\begin{cajaParcial}[title={Ficha t\'ecnica del Parcial 2}]
\renewcommand{\arraystretch}{1.5}
\begin{tabularx}{\textwidth}{>{\bfseries}p{4cm}X}
\textbf{Modalidad:} & Grupal (mismos grupos del TP Integrador), en clase. \\
\textbf{Duraci\'on:} & 60 minutos de redacci\'on + 5 minutos de defensa oral por grupo. \\
\textbf{Cobertura:} & Unidades 3 y 4 + temas transversales (nube, criptograf\'ia, hacking \'etico). \\
\textbf{Entregable:} & Una Pol\'itica de Seguridad Aceptable (PSA) para el caso dado. \\
\textbf{Nota:} & 1 a 10. Se aprueba con 4. \\
\end{tabularx}
\end{cajaParcial}

\subsection{Caso del Parcial 2: ``DevCraft''}

\begin{tcolorbox}[colback=grisClaro, colframe=azulGobierno, arc=3mm]
\textbf{Contexto:} ``DevCraft'' es una empresa de desarrollo de software con 40 empleados. Hace aplicaciones a medida para clientes (bancos, retailers, startups).\\[4pt]
\textbf{Equipo:}
\begin{itemize}[leftmargin=1.5em]
    \item 25 desarrolladores (muchos trabajan remoto o h\'ibrido, usan notebooks propias y corporativas).
    \item 5 dise\~nadores UX/UI.
    \item 3 project managers.
    \item 2 personas en RRHH.
    \item 2 en administraci\'on/finanzas.
    \item 1 gerente general.
    \item 2 personas en ``infraestructura/DevOps'' (manejan la nube de la empresa).
\end{itemize}
\textbf{Situaci\'on:} La empresa creci\'o r\'apido y nunca tuvo una pol\'itica formal. Recientemente un desarrollador subi\'o c\'odigo con credenciales de AWS a un repositorio p\'ublico de GitHub, y un cliente se quej\'o de que su inform\'on confidencial podr\'ia haber quedado expuesta. El director decidi\'o que \textbf{es hora de tener una PSA formal}.
\end{tcolorbox}

\subsection{¿Qu\'e deben entregar?}

Cada grupo debe redactar una \textbf{PSA completa para DevCraft}, con las 9 secciones vistas en la Secci\'on 3, adaptadas al caso. Debe ser:
\begin{itemize}[leftmargin=2em]
    \item \textbf{Realista} para una empresa de 40 personas (no pedir controles de multinacional).
    \item \textbf{Espec\'ifica} sobre los riesgos reales del caso: trabajo remoto/h\'ibrido, BYOD, c\'odigo con credenciales, datos de clientes bancarios.
    \item \textbf{Coherente} con el marco normativo (Ley 25.326, derechos ARCO) y los temas transversales vistos (nube, criptograf\'ia, hacking \'etico).
\end{itemize}

\newpage
\subsection{R\'ubrica oficial del Parcial 2 (sobre 10 puntos)}

\renewcommand{\arraystretch}{1.5}
\begin{tabularx}{\textwidth}{|X|>{\centering\arraybackslash}p{3cm}|}
\hline
\rowcolor{violetaParcial!20}
\textbf{¿Qu\'e se eval\'ua?} & \textbf{Puntos} \\
\hline
Completitud y pertinencia de la PSA para el caso dado & hasta 3 \\
\hline
Alineaci\'on con el marco normativo y de gobierno visto en clase & hasta 2 \\
\hline
Calidad de la defensa oral grupal (claridad, argumentaci\'on) & hasta 2 \\
\hline
Calidad de las respuestas individuales a las preguntas del docente & hasta 3 \\
\hline
\rowcolor{grisClaro}
\textbf{TOTAL} & \textbf{10} \\
\hline
\end{tabularx}

\begin{cajaNota}[title={Atenci\'on: las preguntas individuales valen 30\%}]
El Parcial 2 no solo eval\'ua el documento escrito, sino la \textbf{comprensi\'on individual}. El docente har\'a preguntas cruzadas a cada integrante del grupo. Todos deben poder explicar cualquier parte de la PSA.
\end{cajaNota}

\newpage
% ============================================================
%  SECCIÓN 6 — PLANTILLA DE REDACCIÓN
% ============================================================
\section{Plantilla de redacci\'on para el Parcial 2}

\begin{cajaParcial}[title={Instrucciones}]
Durante el Parcial 2 van a redactar la PSA de DevCraft en esta plantilla. Pueden usar el reverso o hojas adicionales si es necesario. Al terminar, la firman como grupo y la entregan al docente.
\end{cajaParcial}

\vspace{6pt}

\begin{center}
\Large\bfseries\color{azulGobierno}
POL\'ITICA DE SEGURIDAD ACEPTABLE\\[4pt]
\normalsize DevCraft --- Versi\'on 1.0
\end{center}

\vspace{8pt}
\renewcommand{\arraystretch}{1.6}

\noindent\textbf{1. Prop\'osito y alcance}
\begin{tabularx}{\textwidth}{|X|}
\hline
\\[2cm]
\hline
\end{tabularx}

\vspace{6pt}
\noindent\textbf{2. Uso aceptable de los recursos}
\begin{tabularx}{\textwidth}{|X|}
\hline
\\[2cm]
\hline
\end{tabularx}

\vspace{6pt}
\noindent\textbf{3. Contrase\~nas y accesos}
\begin{tabularx}{\textwidth}{|X|}
\hline
\\[2cm]
\hline
\end{tabularx}

\newpage
\noindent\textbf{4. Correo e internet}
\begin{tabularx}{\textwidth}{|X|}
\hline
\\[2cm]
\hline
\end{tabularx}

\vspace{6pt}
\noindent\textbf{5. Dispositivos (incluye BYOD y trabajo remoto)}
\begin{tabularx}{\textwidth}{|X|}
\hline
\\[2.5cm]
\hline
\end{tabularx}

\vspace{6pt}
\noindent\textbf{6. Informaci\'on y datos confidenciales (incluye c\'odigo fuente)}
\begin{tabularx}{\textwidth}{|X|}
\hline
\\[2.5cm]
\hline
\end{tabularx}

\newpage
\noindent\textbf{7. Reporte de incidentes}
\begin{tabularx}{\textwidth}{|X|}
\hline
\\[2cm]
\hline
\end{tabularx}

\vspace{6pt}
\noindent\textbf{8. Cumplimiento y sanciones}
\begin{tabularx}{\textwidth}{|X|}
\hline
\\[2cm]
\hline
\end{tabularx}

\vspace{6pt}
\noindent\textbf{9. Vigencia y firmas}
\begin{tabularx}{\textwidth}{|X|}
\hline
\\[1.5cm]
\hline
\end{tabularx}

\newpage
% ============================================================
%  SECCIÓN 7 — PREPARACIÓN DEFENSA ORAL
% ============================================================
\section{Preparaci\'on para la defensa oral}

Al terminar la redacci\'on, cada grupo tiene \textbf{5 minutos} para presentar su PSA al resto de la clase y al docente. Luego, el docente har\'a preguntas individuales.

\subsection{Estructura recomendada de la presentaci\'on (5 min)}
\begin{enumerate}[leftmargin=2em]
    \item \textbf{1 minuto} — Contexto: por qu\'e DevCraft necesitaba una PSA (el incidente del repositorio).
    \item \textbf{2 minutos} — Los 3-4 puntos m\'as importantes de su PSA (los que priorizaron).
    \item \textbf{1 minuto} — C\'omo van a lograr que los empleados la conozcan y la cumplan (capacitaci\'on, firma al ingreso, recordatorios).
    \item \textbf{1 minuto} — C\'omo se revisar\'a y actualizar\'a la pol\'itica en el tiempo.
\end{enumerate}

\subsection{Preguntas individuales t\'ipicas del docente}

Estas son \textbf{ejemplos reales} de preguntas que el docente puede hacer a cada integrante del grupo. Pract\'iquenlas antes de la defensa:

\renewcommand{\arraystretch}{1.5}
\begin{tabularx}{\textwidth}{|X|}
\hline
\textbf{Sobre el caso:}\\[2pt]
--- ¿Por qu\'e incluyeron esta regla sobre c\'odigo fuente y repositorios?\\
--- ¿C\'omo manejar\'ian el hecho de que muchos desarrolladores trabajan con notebooks propias?\\
--- ¿Qu\'e hacen si un empleado se niega a firmar la PSA?\\
\hline
\textbf{Sobre el marco normativo:}\\[2pt]
--- ¿Su PSA respeta la Ley 25.326? ¿C\'omo?\\
--- ¿Qu\'e derechos ARCO deber\'ia conocer cada empleado?\\
--- ¿C\'omo se relaciona esta PSA con la pol\'itica de privacidad de la empresa?\\
\hline
\textbf{Sobre temas transversales:}\\[2pt]
--- Si DevCraft usa AWS, ¿qu\'e parte de la seguridad es responsabilidad de la empresa y cu\'al de AWS?\\
--- ¿C\'omo se conecta esta PSA con lo que vieron de hashing y criptograf\'ia?\\
--- ¿Qu\'e rol tiene la concientizaci\'on (clase 10) para que la PSA se cumpla?\\
\hline
\textbf{Sobre gobierno:}\\[2pt]
--- ¿Qui\'en es el responsable \'ultimo de la seguridad en DevCraft?\\
--- ¿C\'omo se revisa y actualiza esta pol\'itica?\\
--- ¿Qu\'e pasa si un gerente incumple la pol\'itica?\\
\hline
\end{tabularx}

\begin{cajaNota}[title={Consejos para la defensa}]
\begin{enumerate}[leftmargin=2em]
    \item \textbf{Distribuyan roles} para la presentaci\'on: qui\'en introduce, qui\'en explica los puntos clave, qui\'en cierra.
    \item \textbf{Todos deben estar listos para responder} preguntas individuales. Repasen la PSA juntos antes de presentar.
    \item \textbf{Sean honestos} si no saben algo: es mejor decir ``lo discutimos pero no lo resolvimos'' que inventar.
    \item \textbf{Usen ejemplos concretos} del caso DevCraft, no generalidades.
\end{enumerate}
\end{cajaNota}

\newpage
% ============================================================
%  SECCIÓN 8 — GLOSARIO Y CIERRE
% ============================================================
\section{Glosario de la Clase 11}

\renewcommand{\arraystretch}{1.4}
\begin{tabularx}{\textwidth}{>{\bfseries}p{4.5cm}X}
\toprule
\textbf{T\'ermino} & \textbf{Definici\'on en una l\'inea} \\
\midrule
Gobierno de la seguridad & Estructura de pol\'iticas, roles y procesos que dirigen la seguridad. \\
Pol\'itica & Declaraci\'on de intenci\'on de la direcci\'on (el ``qu\'e'' y el ``por qu\'e''). \\
Norma & Regla espec\'ifica y obligatoria que desarrolla la pol\'itica. \\
Procedimiento & Pasos operativos detallados para cumplir una norma. \\
PSA / AUP & Pol\'itica de Seguridad Aceptable (Acceptable Use Policy). \\
CISO & Director de Seguridad de la Informaci\'on (rol ejecutivo). \\
DPO & Delegado de Protecci\'on de Datos. \\
Due\~no del activo & Persona responsable de un activo espec\'ifico. \\
BYOD & ``Bring Your Own Device''; empleados usan dispositivos propios. \\
Vigencia & Fecha desde la cual una pol\'itica es v\'alida. \\
\bottomrule
\end{tabularx}

\vspace{1cm}

\section{Cierre y puente hacia la Clase 12}

\begin{cajaGobierno}[title={Lo que nos llevamos de la Clase 11}]
\begin{itemize}[leftmargin=2em]
    \item La diferencia entre \textbf{seguridad t\'ecnica} y \textbf{gobierno de la seguridad}.
    \item La \textbf{pir\'amide documental}: pol\'iticas $\rightarrow$ normas $\rightarrow$ procedimientos.
    \item La estructura de una \textbf{PSA} completa y su importancia legal y cultural.
    \item Los \textbf{roles} clave en seguridad (CISO, DPO, due\~nos de activos).
    \item Experiencia real de redactar y defender una PSA (Parcial 2).
\end{itemize}
\end{cajaGobierno}

\begin{cajaNota}[title={Puente hacia la Clase 12: Factor humano e incidentes}]
Una PSA es el \textbf{piso} del gobierno. Pero ni la mejor pol\'itica evita que los incidentes ocurran. La pr\'oxima clase vamos a ver qu\'e hacer \textbf{cuando algo sale mal}: c\'omo gestionar un incidente, c\'omo reacciona la organizaci\'on, y por qu\'e el factor humano es a la vez el eslab\'on m\'as d\'ebil y el m\'as importante.\\[4pt]
Adem\'as, prepararemos la \textbf{Junta de Crisis} de la Clase 12 (un juego de roles donde cada integrante toma un papel durante una filtraci\'on de datos).
\end{cajaNota}

\vspace{1cm}

\section{Espacio para notas y preparaci\'on del Parcial 2}

\begin{tcolorbox}[colback=grisClaro, colframe=gray!50, arc=2mm, height=7cm]
\textcolor{gray}{\itshape Us\'a este espacio para anotar ideas para la PSA, dudas sobre el caso DevCraft, o posibles respuestas a las preguntas individuales del docente.}
\end{tcolorbox}

\vfill
\begin{center}
\small\textcolor{gray}{%
Cuadernillo elaborado para la c\'atedra de Seguridad de la Informaci\'on y de
Sistemas --- Lic.\ en Gesti\'on de Negocios Digitales, UCA.\\
Prohibida su distribuci\'on fuera del \'ambito de la cursada sin autorizaci\'on del docente.}
\end{center}

\end{document}